\section{\texttt{boot.asm} 详解：从物理地址到高半区}

\texttt{boot.asm} 是内核的入口点 \texttt{\_start}。
它的核心任务是：在没有分页的情况下构建临时页目录，开启分页后跳转到高半区。

\subsection{保存 Multiboot2 信息}

\codefile{src/boot/boot.asm}
\begin{lstlisting}[style=nasmstyle]
_start:
    mov esi, eax            ; magic (0x36d76289)
    mov edi, ebx            ; mb2 info pointer (physical address)
\end{lstlisting}

GRUB 传入的 \reg{EAX}（magic）和 \reg{EBX}（info 指针）保存在 callee-saved
寄存器 \reg{ESI}/\reg{EDI} 中，以便后续传给 \texttt{kmain}。

\subsection{PSE 页目录构造}

\codefile{src/boot/boot.asm}
\begin{lstlisting}[style=nasmstyle]
; 清零页目录 (1024 x 4B = 4 KiB)
push edi
mov ecx, 1024
xor eax, eax
mov edi, boot_pd
rep stosd
pop edi

; 填充 PSE 映射: 256 x 4 MiB = 1 GiB
xor ecx, ecx
.fill_pse:
    mov eax, ecx
    shl eax, 22             ; eax = ecx * 4 MiB
    or  eax, 0x83           ; PS(bit7) | RW(bit1) | Present(bit0)
    mov [boot_pd + ecx*4], eax          ; identity: PD[i]
    mov [boot_pd + 768*4 + ecx*4], eax  ; high-half: PD[768+i]
    inc ecx
    cmp ecx, 256
    jb  .fill_pse
\end{lstlisting}

\begin{whybox}
为什么用 PSE（4\,MiB 大页）而不是标准 4\,KiB 页？

4\,KiB 页需要分配 page table（每个 PDE 指向一个 4\,KiB 页表），
但此时 PMM 还没初始化——我们甚至不知道哪些物理内存可用。
PSE 模式下 PDE 自身就包含 4\,MiB 页的物理基地址，不需要页表。
只需一个页目录（4\,KiB）就能映射整个地址空间。

\texttt{vmm\_init()} 之后会用 4\,KiB 页表替换这些 PSE 条目，
实现更细粒度的权限控制。
\end{whybox}

PSE PDE 的位格式：

\begin{figure}[H]
  \centering
  % figures/ch02/fig02-pse-pde-bits.tex — auto-extracted from chapter source
\begin{tikzpicture}[
  bitcell/.style={draw, minimum width=0.7cm, minimum height=0.6cm,
                  font=\ttfamily\small, anchor=west},
]
  \node[bitcell, minimum width=4.2cm, fill=blue!12] at (0, 0) {\small phys[31:22]};
  \node[bitcell, minimum width=2.8cm, fill=gray!10] at (4.2, 0) {\small 保留/PAT};
  \node[bitcell, fill=red!15]   at (7, 0) {\small 1};
  \node[bitcell, fill=gray!10]  at (7.7, 0) {\small 0};
  \node[bitcell, fill=gray!10]  at (8.4, 0) {\small 0};
  \node[bitcell, fill=gray!10]  at (9.1, 0) {\small 0};
  \node[bitcell, fill=gray!10]  at (9.8, 0) {\small 0};
  \node[bitcell, fill=green!15] at (10.5, 0) {\small 1};
  \node[bitcell, fill=green!15] at (11.2, 0) {\small 1};

  \node[font=\scriptsize, above] at (2.1, 0.6) {bits 31--22};
  \node[font=\scriptsize, above] at (7.0, 0.6) {PS};
  \node[font=\scriptsize, above] at (10.5, 0.6) {RW};
  \node[font=\scriptsize, above] at (11.2, 0.6) {P};

  \node[font=\scriptsize, below] at (6, -0.3) {0x83 = 1000\,0011$_2$ = PS $|$ RW $|$ Present};
\end{tikzpicture}

  \caption{PSE 模式 PDE 位域（\texttt{0x83} 标志）}
  \label{fig:pse-pde-bits}
\end{figure}

\subsection{双重映射：Identity + High-Half}
\label{sec:dual-mapping}

循环将同一物理地址同时写入两个 PD 位置：

\begin{itemize}
  \item \texttt{PD[i]}（$i = 0..255$）——Identity 映射：虚拟 $i \times 4\text{MiB}$ → 物理 $i \times 4\text{MiB}$
  \item \texttt{PD[768+i]}——High-half 映射：虚拟 $\hex{C0000000} + i \times 4\text{MiB}$ → 物理 $i \times 4\text{MiB}$
\end{itemize}

\begin{figure}[H]
  \centering
  % figures/ch02/fig03-dual-mapping.tex — auto-extracted from chapter source
\begin{tikzpicture}[
  block/.style={draw, minimum width=4.5cm, minimum height=0.7cm,
                font=\small, anchor=west},
  arr/.style={-{Stealth[length=5pt]}, thick, blue!60!black},
  addr/.style={font=\ttfamily\scriptsize\color{gray}, anchor=east},
]
  % 虚拟地址空间
  \node[font=\bfseries\small] at (1.5, 3) {虚拟地址空间};
  \node[block, fill=green!15] (id) at (0, 2) {identity: 0--1\,GiB};
  \node[block, fill=gray!10]  (gap) at (0, 1.3) {未映射};
  \node[block, fill=red!15]   (hh) at (0, 0.6) {high-half: 0xC0000000--};
  \node[addr] at (-0.2, 2) {\hex{00000000}};
  \node[addr] at (-0.2, 1.3) {\hex{40000000}};
  \node[addr] at (-0.2, 0.6) {\hex{C0000000}};

  % 物理地址空间
  \node[font=\bfseries\small] at (9, 3) {物理内存};
  \node[block, fill=blue!12, minimum height=1.4cm] (phys) at (7, 1.3)
    {RAM: 0--1\,GiB};
  \node[addr] at (6.8, 2) {\hex{00000000}};

  % 箭头
  \draw[arr] (4.5, 2.0) -- (7, 2.0) node[midway, above, font=\scriptsize] {1:1};
  \draw[arr] (4.5, 0.6) -- (7, 1.0) node[midway, below, font=\scriptsize]
    {-\,\hex{C0000000}};
\end{tikzpicture}

  \caption{引导期双重映射：identity 和 high-half 指向同一物理内存}
  \label{fig:dual-mapping}
\end{figure}

\begin{warnbox}
如果只建 high-half 映射而不建 identity 映射会发生什么？

开启分页后，CPU 的下一条指令在物理地址 $\sim$\hex{001xxxxx}。
没有 identity 映射，MMU 无法翻译当前 EIP → \textbf{Page Fault} →
此时 IDT 还没设置 → \textbf{Double Fault} →
没有 DF handler → \textbf{Triple Fault} → CPU 重启。

整个"死亡连锁"发生在一条指令之内，QEMU 会直接 reboot，
串口什么输出都看不到——这也是引导调试最令人头疼的问题之一。
\end{warnbox}

\subsection{开启 PSE 和分页}

\codefile{src/boot/boot.asm}
\begin{lstlisting}[style=nasmstyle]
; CR4.PSE = 1 (允许 4 MiB 页)
mov eax, cr4
or  eax, 0x00000010     ; bit 4 = PSE
mov cr4, eax

; CR3 <- boot_pd (物理地址)
mov eax, boot_pd
mov cr3, eax

; CR0.PG = 1 (开启分页)
mov eax, cr0
or  eax, 0x80000000     ; bit 31 = PG
mov cr0, eax
\end{lstlisting}

\begin{cpustate}
三步 CR 操作后 CPU 状态的逐步变化：
\begin{enumerate}
  \item \texttt{CR4.PSE = 1}——CPU 现在\emph{允许} 4\,MiB 页，但还没开启分页。
  \item \texttt{CR3 = boot\_pd}——告诉 CPU 页目录在哪，但分页仍关闭。
  \item \texttt{CR0.PG = 1}——分页\textbf{立即生效}！
        从下一条指令开始，所有内存访问经过 MMU 翻译。
        Identity 映射保证当前 EIP 仍然有效。
\end{enumerate}
\end{cpustate}

\subsection{跳转到高半区}

\codefile{src/boot/boot.asm}
\begin{lstlisting}[style=nasmstyle]
mov eax, _start_high    ; 链接地址 = 0xC01xxxxx
jmp eax                 ; 跳到高半区虚拟地址
\end{lstlisting}

\texttt{\_start\_high} 在 \texttt{.text} 段中，链接在 \hex{C01xxxxx}。
\texttt{mov eax, \_start\_high} 把这个高虚拟地址加载到 \reg{EAX}，
\texttt{jmp eax} 跳转——MMU 通过 PD[768] 将 \hex{C01xxxxx} 翻译回物理 \hex{001xxxxx}。

从此内核运行在高半区虚拟地址空间。

% ============================================================================

